
\subsection{Causal Reasoning}
\label{subsec:p3_s7_1_6}
% - Cause--effect relationships between spatial events (``why'' and ``what-if'').
% - Structural Causal Models (SCM), Counterfactual SCM, PC Algorithm, GES, LiNGAM
% - Task dependency planning, event explanation, intervention & prevention, multi-robot coordination

Causal reasoning moves from describing \emph{what} happened to explaining \emph{why} it happened and predicting \emph{what would have happened} under a different action. For a robot the relevant outcomes are chains of spatial events (a collision, a delay, a deadlock, an unstable grasp, a near-miss in shared space), and making the causal assumptions explicit, rather than leaving them implicit in ``actions cause state changes'', is what lets the system attribute a failure to its causes rather than its symptoms, evaluate counterfactual alternatives such as a different approach direction or timing, and, with several agents, separate primary causes from downstream effects when choosing an intervention. Figure~\ref{fig:p3_s7_1_6_causal_reasoning_abstract} gives the abstract formulation.

\begin{figure}[!th]
    \centering
    \includegraphics[width=0.99\linewidth]{figures/part3/section7/716_causal.pdf}
    \caption{Abstract task formulation of causal reasoning for spatial intelligence: (i) event streams, candidate variables, prior knowledge, and interventional data feed (ii) methods that discover and verify causal structure, yielding (iii) causal models, explanations, counterfactual answers, and intervention or prevention strategies. The method families serve different roles rather than forming a serial pipeline. Verification returns a satisfying witness or an unsatisfiability proof/core where supported; counterfactual scenes keep the objects fixed while changing the blue agent's action.}
    \label{fig:p3_s7_1_6_causal_reasoning_abstract}
\end{figure}

\subsubsection{Task Definition}
The inputs are spatiotemporal event streams (agent actions, object states, contact and visibility relations, environment variables) mapped to causal variables by measurement functions (\texttt{in\_contact}, \texttt{blocked}, \texttt{priority\_violated}), optional prior knowledge from physics or commonsense, and, when available, interventional data in which a controllable variable was set deliberately. A structural causal model over these variables,
\begin{equation}
    \mathcal{M} = \langle \mathcal{G}, \{f_i\}, P(U) \rangle, \qquad X_i := f_i(\mathrm{Pa}(X_i), U_i),
\end{equation}
consists of a directed acyclic graph $\mathcal{G}$, a mechanism $f_i$ that generates each variable from its parents $\mathrm{Pa}(X_i)$ and exogenous noise $U_i$, and a noise distribution $P(U)$~\cite{Pearl2009Causality}. The outputs are the model itself when it must be learned, the explanation of an observed outcome (which variables caused it), the answer to an interventional query $P(Y \mid do(X=x))$ or a counterfactual one ($Y_{x}$ for a unit whose actual trace was observed), and an intervention or prevention strategy that changes controllable causes to reach a desired spatial outcome.

\subsubsection{Methodology}
The methods differ in where the causal structure comes from (assumed, learned from data, or written as domain knowledge) and in what they are asked (intervention effects, counterfactuals, responsibility, verification); in spatial domains they also differ in how they use locality, temporal precedence, and network connectivity as inductive bias.

\paragraph{Structural Causal Models and Counterfactuals.}
An SCM makes mechanisms explicit (``clearance determines collision risk'', ``traffic density sets delay'') and separates intervention from observation: the $do(\cdot)$ operator replaces a structural equation instead of conditioning on a value, which is what is needed when, say, congestion and delay are both driven by time of day. Spatial pipelines typically start from a graph fixed by temporal precedence and locality priors (edges among nearby regions, interacting agents, adjacent cells), estimate the mechanisms from data, and validate them against implied conditional independences and refutation tests~\cite{Peters2017Elements}. Counterfactuals follow the abduction-action-prediction recipe: infer the exogenous noise consistent with the observed trace, apply the hypothetical intervention, predict the outcome~\cite{Pearl2009Causality}. DoWhy and DoWhy-GCM~\cite{DoWhy2020,Bloebaum2024DoWhyGCM} implement identification, estimation, and refutation for this workflow on real spatiotemporal data.

\paragraph{Causal Discovery: Constraint- and Score-Based.}
When the graph is unknown but observational data are plentiful, constraint-based methods read structure off conditional independences: the PC algorithm builds the skeleton by independence tests and orients what it can, returning a Markov equivalence class, and FCI~\cite{Spirtes2000CPS} handles latent confounding and selection bias by returning a partial ancestral graph. Score-based methods instead search graph space under an objective such as BIC: GES~\cite{Chickering2002GES} greedily adds and then removes edges over equivalence classes, and its forward-backward structure often recovers mesoscale organization such as bottlenecks that drive downstream congestion. Spatial data strain both. Large conditioning sets are expensive and statistically weak, so practice restricts conditioning to a radius or $k$-hop neighborhood, forbids edges that violate adjacency or reachability priors, and uses order-independent variants (PC-stable~\cite{Colombo2014PCStable}) whose high-dimensional consistency is understood~\cite{Kalisch2007PC}; ambiguous orientations are then resolved by targeted interventions. causal-learn~\cite{zheng2024causal} and the TETRAD / Py-Tetrad implementations of FGES~\cite{Ramsey2018Tetrad,Ramsey2023PyTetrad} provide these pipelines at facility and city scale.

\paragraph{Independent-Noise Discovery (LiNGAM).}
In linear-Gaussian settings many edge directions are unidentifiable from independences alone. LiNGAM~\cite{Shimizu2006LiNGAM} assumes a linear structural model with independent non-Gaussian noise, under which the causal ordering is identifiable, and DirectLiNGAM~\cite{Shimizu2011DirectLiNGAM} estimates it by regression rather than ICA. Spatiotemporal signals need care because autocorrelation violates i.i.d. sampling assumptions and aggregation may change the noise distribution; temporal blocking, prewhitening, and locality constraints are the usual remedies, and the recovered ordering then scaffolds an SCM for intervention analysis.

\paragraph{Causal Action Languages and Plan-Space Causal Links.}
Much causal structure in robotics is designed knowledge rather than discovered: actions have direct and indirect effects, and the system must reason about change under incomplete information. The action language C\textsuperscript{+}~\cite{Giunchiglia2004CPlus} writes causal theories of action with inertia and nonmonotonic semantics, compactly describing how fluents such as \texttt{occupied(c)} or \texttt{visible(a,o)} evolve; CCALC~\cite{Lifschitz2000CCalc} compiles them to SAT so that prediction, diagnosis, and plan synthesis use the same solvers; and multi-valued fluents and external geometric predicates (collision checks, IK) are being integrated without abandoning the formal semantics~\cite{giunchiglia2023causal}. In plan space, partial-order causal-link planning~\cite{McAllester1991SNLP} represents a plan as partially ordered actions plus links $A \xrightarrow{p} B$ (action $A$ establishes precondition $p$ for $B$), which makes threats explicit and repairable by promotion, demotion, or binding constraints; UCPOP~\cite{Penberthy1992UCPOP} extended the representation to ADL and conditional effects and VHPOP~\cite{Younes2004VHPOP} added heuristics and reachability analysis. In spatial domains a causal link is often a physical dependency, ``unblock door'' enabling ``move cart'', which bridges symbolic causality and spatial feasibility.

\paragraph{Execution and Verification.}
A causal plan must be realized in time and checked. Simple temporal networks~\cite{Dechter1991STN} enforce difference constraints between events with polynomial consistency and dispatchable forms~\cite{Tsamardinos1998Dispatch}, and STNUs with dynamic-controllability analysis~\cite{Morris2001STNU} handle contingent durations, so that ordering and clearance constraints stay satisfiable under bounded disturbance. The SAT/SMT machinery of §\ref{subsec:p3_s7_1_4} enters here as an explanation device: when a schedule is unsatisfiable, the unsat core is a set of conflicting commitments, a concrete statement of which decisions caused the failure~\cite{Surynek2016SATMAPF,Yu2016OptimalMAPF}; and for hybrid dynamics $\delta$-complete SMT (dReal)~\cite{Gao2013dReal} and dReach~\cite{Kong2015dReach} verify safety and synthesize interventions with numerically sound proofs or counterexamples.

\subsubsection{Benchmarks}
No single benchmark covers the causal pipeline, and the suites stress different links of it. CLEVRER~\cite{Yi2020CLEVRER} pairs synthetic collision videos with descriptive, explanatory, predictive, and counterfactual questions, and CLEVRER-Humans~\cite{CLEVRERHumans2024} adds human-annotated causal event graphs against which model attributions can be compared; the video suites of §\ref{subsec:p3_s7_1_2} (PHYRE~\cite{bakhtin2019phyre}, Physion and Physion++~\cite{Bear2021Physion,tung2023physion}) score intervention selection and latent-property inference, and CRaFT~\cite{Saha2022CRaFT}, ContPhy~\cite{Zheng2024ContPhy}, NExT-QA~\cite{Xiao2021NExTQA}, AGQA~\cite{GrundeMcLaughlin2021AGQA,Gupta2023AGQA2}, and SUTD-TrafficQA~\cite{Xu2021TrafficQA} ask causal questions in natural language about synthetic and real-world videos, including traffic scenes. Embodied platforms with predicate-level state (BEHAVIOR-1K, ALFRED, iGibson 2.0~\cite{Li2021iGibson2}, ManiSkill2~\cite{Gu2023ManiSkill2}, RLBench~\cite{james2020rlbench}; §\ref{subsec:p3_s7_1_5}) test whether inferred dependencies improve intervention selection and recovery; the MAPF suites and Flatland~\cite{Mohanty2020Flatland} (§\ref{subsec:p3_s7_1_3}) test reasoning over space-time constraints in which one agent's decision causes another's delay; and ARCH-COMP~\cite{ARCHCOMP2024} with the dReal / dReach problem sets~\cite{Gao2013dReal,Kong2015dReach} evaluate causal safety reasoning on hybrid dynamics by proof or counterexample and solve time. Discovery itself is evaluated on recovery of the Markov equivalence class under faithfulness, with scalability bounded by conditioning-set size.

\subsubsection{Challenges and Trends}
Identification is brittle. Latent confounding, selection bias, and temporal aggregation can make a query non-identifiable from logs, and estimates become unstable when locality and precedence assumptions fail~\cite{Pearl2009Causality,Peters2017Elements}; constraint- and score-based discovery is expensive on dense telemetry with strong short-range dependence and sensitive to test strength, ordering, and regularization~\cite{Spirtes2000CPS,Colombo2014PCStable,Kalisch2007PC,Chickering2002GES}; and LiNGAM's non-Gaussianity is exactly what smoothed trajectories and aggregated traffic signals lack~\cite{Shimizu2006LiNGAM,Shimizu2011DirectLiNGAM}. The response is structure-aware discovery that constrains search by topology and locality and then resolves the remaining ambiguity by targeted robot interventions, connecting learning to actuation~\cite{Spirtes2000CPS,Chickering2002GES}.

Causal reasoning is only as good as its variables, and in embodied settings contact, occlusion, blockage, and object state must be grounded from noisy perception. The errors of that mapping propagate into SCM estimates and counterfactual conclusions~\cite{Peters2017Elements,DoWhy2020}. Knowledge-enabled frameworks such as KnowRob and CRAM (§\ref{subsec:p3_s7_1_5}) show the shape of the interface, perception populating an explicit, interpretable model whose parameters are then estimated statistically, and the trend is toward such perception-to-symbol interfaces that tolerate sensor noise~\cite{tenorth2013knowrob,beetz2010cram,beetz2023cram}.

Between a causal recommendation and a safe execution there is a verification-synthesis gap: the recommendation must become a temporally controllable schedule and a certified controller, and execution uncertainty invalidates a symbolic causal plan unless temporal networks and hybrid verification are integrated with it~\cite{Dechter1991STN,Morris2001STNU}. SCM toolchains are therefore being embedded in decision stacks so that counterfactual policies compile into action-language or plan-space structures and then into dispatchable networks~\cite{Bloebaum2024DoWhyGCM,Giunchiglia2004CPlus,Morris2001STNU}, with SAT/SMT certificates and unsat-core explanations supplying the accountability~\cite{Surynek2016SATMAPF,Gao2013dReal,Kong2015dReach}; evaluation that scores explanation quality, interventional effectiveness, and certification cost on the same task is what the field still lacks.
